\section{Conclusion}


\subsection{RQ 1: Terrain generation}
The system we have created currently supports 3 features: The surface terrain is generated using Fractal Perlin noise as a two-dimensional heightmap, and split into separate chunks. For each height map, a mesh is generated. This method is able to generate simple terrain with some ruggedness, which is enough for simple environments, such as forests.

To generate rivers, we randomly plan the path of a meandering river, which randomly turns at each step of the river planning. This path is transformed into a fat curve which is used in other steps of the program. Using the cross section of the river, the fat curve modifies the height map of the terrain to imprint the river bed and creates a mesh that represents the river surface. Although the generation of the river itself works, the river planning is very simplistic and often generates rivers that are not visible from the camera's perspective. From our experiments more than half of the terrains we have generated does not include a significant portion of the river in the view.

For the generation of vegetation we have selected two methods. The first method randomly places objects on a Poisson Disk Distribution, where all objects have a minimum distance where no other object can appear. With the other approach, a probability distribution is used to randomly place objects with each object modifying the distribution. Using different kernels changes the behaviour of objects placed, where certain kernels produce a prohibiting effect for other objects, while other kernels can attract objects to clutter near each other. Using both types of generation, different types of behaviour can be produced when generating the environment. To prevent overlap with rivers, the fat curve of the river acts as a filter for PDD objects, where objects are not generated if they are placed on a river. For the deformation kernels, the fat curve modifies the initial probability distribution to change the probability of an object placed on a certain part of the river. To place the objects onto the terrain, the height map is used to find the vertical position of the object on the terrain and the update is placed on the surface.

\subsection{RQ 2: Parameters of the system}
In the architecture we have created, each of the different components has a number of parameters that influence the generation or transformation of the specific feature. These parameters are categorised into two categories, where parameters can be considered as constant in most of the generated terrains, and parameters that change the terrain in order to get a desired effect.

For generating the terrain, most parameters are connected to the mathematical noise functions, and the influence of those parameters requires knowledge of these noise functions. The generation of rivers has parameters that describe the general shape and direction of the river. For the generation of objects, the parameters are inherent to the objects to be placed on the terrain. In the Poisson Disk Distribution, objects specify separate radii for objects of the same type or objects in a lower layer. In the deformation field, the objects placed specify how the probability field is modified in an area around the object. The amount of objects placed using the PDD is limited by the bounds of the terrain and the minimum distance, while the deformation uses a parameter to control the amount of objects generated.

\subsection{RQ 3: Validating terrain}
To validate the generated environments, we have selected three qualities of the terrain. For the first metric, we analyse the density of the generated trees. With the PDD algorithm, a minimum distance between each object is ensured, but by increasing this minimum distance, the density of objects decreases. This relation is an inverse quadratic relation and can allow the designer to estimate a density of objects the PDD algorithm will generate. When a certain density of objects is required, it can be calculated back to a minimum radius and then validated when the terrain is generated by counting each object.

For the second quality, the visibility of the generated rivers is analysed. In this metric, the generated terrain is evaluated by the proportion of river visible on the terrain. This method only analyses the surface, and therefore placed objects are removed. From this terrain object segmentation is used to calculate the amount of pixels that are either river or terrain. From this a ratio is calculated between the surface and the rivers, that can be used by a designer to find an acceptable range where a terrain contains enough visible river. When using multiple camera positions, such as each side of the terrain as is used as a method to improve the visibility of rivers, the same terrain can be used in some cases to provide a better view of the river, although this is not the case for all terrains.

For the placement objects, we analysed the visibility of each object, and looked at whether some objects take up a significant portion of the screen. Instead of object segmentation, panoptic segmentation is used where each individial object is separated. For each terrain, the pixels per objects are counted and calculated as the proportion of the total screen size. From this we have seen that some objects occupy a significant portion of the screen. With these proportions, large objects can be removed if they block a large part of the view, or find a minimum distance where most large objects appear in and remove all objects within this distance.
